\documentclass[journal,onecolumn,11pt]{IEEEtran}
% Supplementary material for the DistilSCD manuscript (IEEE JSTARS).
\usepackage[utf8]{inputenc}
\usepackage[T1]{fontenc}
\usepackage{amsmath,amssymb}
\usepackage{graphicx}
\usepackage{booktabs,array}
\usepackage{siunitx}
\usepackage[hidelinks]{hyperref}
\usepackage{cite}
\usepackage{url}
% Generated by paper/scripts/make_numbers.py from paper/tables/*.csv - do not edit by hand.
\newcommand{\SecBaseN}{7}
\newcommand{\SecBaseSek}{23.90}
\newcommand{\SecBaseSekSd}{0.15}
\newcommand{\SecBaseSekFinal}{23.65}
\newcommand{\SecBaseSekFinalSd}{0.22}
\newcommand{\SecBaseFscd}{64.42}
\newcommand{\SecBaseMiou}{73.23}
\newcommand{\SecBaseSekFinalTwoHundred}{23.30}
\newcommand{\SecChgN}{5}
\newcommand{\SecChgSek}{25.59}
\newcommand{\SecChgSekSd}{0.13}
\newcommand{\SecChgSekFinal}{25.57}
\newcommand{\SecChgSekFinalSd}{0.13}
\newcommand{\SecChgFscd}{65.62}
\newcommand{\SecChgMiou}{74.37}
\newcommand{\SecChgGain}{1.69}
\newcommand{\SecChgGainLo}{1.51}
\newcommand{\SecChgGainHi}{1.88}
\newcommand{\SecChgGainP}{<10^{-4}}
\newcommand{\SecChgGainFinal}{1.93}
\newcommand{\SecChgGainFinalLo}{1.70}
\newcommand{\SecChgGainFinalHi}{2.15}
\newcommand{\SecChgVsT}{$-$0.52}
\newcommand{\SecChgVsTLo}{$-$0.68}
\newcommand{\SecChgVsTHi}{$-$0.35}
\newcommand{\SecChgVsTP}{=0.0010}
\newcommand{\SecChgVsTFinal}{$-$0.53}
\newcommand{\SecChgVsTFinalP}{=0.0007}
\newcommand{\SecChgGapClosed}{77}
\newcommand{\SecFullN}{5}
\newcommand{\SecFullSek}{26.20}
\newcommand{\SecFullSekSd}{0.11}
\newcommand{\SecFullSekFinal}{26.20}
\newcommand{\SecFullSekFinalSd}{0.11}
\newcommand{\SecFullFscd}{66.29}
\newcommand{\SecFullMiou}{74.57}
\newcommand{\SecFullGain}{2.30}
\newcommand{\SecFullGainLo}{2.13}
\newcommand{\SecFullGainHi}{2.47}
\newcommand{\SecFullGainP}{<10^{-4}}
\newcommand{\SecFullGainFinal}{2.55}
\newcommand{\SecFullGainFinalLo}{2.33}
\newcommand{\SecFullGainFinalHi}{2.77}
\newcommand{\SecFullVsT}{+0.09}
\newcommand{\SecFullVsTLo}{$-$0.04}
\newcommand{\SecFullVsTHi}{+0.23}
\newcommand{\SecFullVsTP}{=0.13}
\newcommand{\SecFullVsTFinal}{+0.09}
\newcommand{\SecFullVsTFinalP}{=0.15}
\newcommand{\SecFullGapClosed}{104}
\newcommand{\SecVmbN}{5}
\newcommand{\SecVmbSek}{25.69}
\newcommand{\SecVmbSekSd}{0.17}
\newcommand{\SecVmbSekFinal}{25.69}
\newcommand{\SecVmbSekFinalSd}{0.18}
\newcommand{\SecVmbFscd}{65.73}
\newcommand{\SecVmbMiou}{74.38}
\newcommand{\SecVmbGain}{1.79}
\newcommand{\SecVmbGainLo}{1.57}
\newcommand{\SecVmbGainHi}{2.01}
\newcommand{\SecVmbGainP}{<10^{-4}}
\newcommand{\SecVmbGainFinal}{2.04}
\newcommand{\SecVmbGainFinalLo}{1.78}
\newcommand{\SecVmbGainFinalHi}{2.30}
\newcommand{\SecVmbVsT}{+0.38}
\newcommand{\SecVmbVsTLo}{+0.16}
\newcommand{\SecVmbVsTHi}{+0.59}
\newcommand{\SecVmbVsTP}{=0.009}
\newcommand{\SecVmbVsTFinal}{+0.37}
\newcommand{\SecVmbVsTFinalP}{=0.009}
\newcommand{\SecVmbGapClosed}{127}
\newcommand{\SecFullTwoHundredSek}{26.32}
\newcommand{\SecFullTwoHundredSekSd}{0.11}
\newcommand{\SecFullTwoHundredSekFinal}{26.27}
\newcommand{\SecFullTwoHundredVsT}{+0.22}
\newcommand{\SecFullTwoHundredVsTLo}{+0.08}
\newcommand{\SecFullTwoHundredVsTHi}{+0.35}
\newcommand{\SecFullTwoHundredVsTP}{=0.01}
\newcommand{\SecFullTwoHundredGain}{2.42}
\newcommand{\SecFullTwoHundredGainFinal}{2.96}
\newcommand{\SemInc}{0.61}
\newcommand{\SemIncLo}{0.42}
\newcommand{\SemIncHi}{0.80}
\newcommand{\SemIncP}{=0.0009}
\newcommand{\SemIncPos}{5/5}
\newcommand{\SemIncMin}{0.45}
\newcommand{\SemIncMax}{0.78}
\newcommand{\TeacherDiff}{0.51}
\newcommand{\TeacherDiffLo}{0.27}
\newcommand{\TeacherDiffHi}{0.76}
\newcommand{\TeacherDiffP}{=0.004}
\newcommand{\TeacherDiffPos}{5/5}
\newcommand{\TeacherDiffMin}{0.24}
\newcommand{\TeacherDiffMax}{0.70}
\newcommand{\VmbVsChg}{0.10}
\newcommand{\VmbVsChgLo}{0.02}
\newcommand{\VmbVsChgHi}{0.18}
\newcommand{\VmbVsChgP}{=0.02}
\newcommand{\VmbVsChgPos}{5/5}
\newcommand{\VmbVsChgMin}{0.05}
\newcommand{\VmbVsChgMax}{0.21}
\newcommand{\ExtInc}{0.12}
\newcommand{\ExtIncLo}{0.05}
\newcommand{\ExtIncHi}{0.19}
\newcommand{\ExtIncP}{=0.008}
\newcommand{\ExtIncPos}{5/5}
\newcommand{\ExtIncMin}{0.06}
\newcommand{\ExtIncMax}{0.18}
\newcommand{\TeacherVitgSek}{26.11}
\newcommand{\TeacherVitgSekExact}{26.1079}
\newcommand{\TeacherVmbSek}{25.31}
\newcommand{\TeacherVmbSekExact}{25.314}
\newcommand{\StudentParams}{16.48}
\newcommand{\TeacherParams}{548.17}
\newcommand{\RatioParams}{33.3}
\newcommand{\StudentGmacs}{29.48}
\newcommand{\TeacherGmacs}{1509.74}
\newcommand{\RatioGmacs}{51.2}
\newcommand{\StudentGmacsFvcore}{31.42}
\newcommand{\StudentLatFpThirtyTwo}{129.4}
\newcommand{\TeacherLatFpThirtyTwo}{1613.4}
\newcommand{\RatioLatFpThirtyTwo}{12.5}
\newcommand{\StudentLatBf}{113.7}
\newcommand{\TeacherLatBf}{561.1}
\newcommand{\RatioLatBf}{4.9}
\newcommand{\StudentLatFpSixteen}{112.6}
\newcommand{\TeacherLatFpSixteen}{487.1}
\newcommand{\RatioLatFpSixteen}{4.3}
\newcommand{\StudentLatOneFpThirtyTwo}{26.3}
\newcommand{\TeacherLatOneFpThirtyTwo}{268.5}
\newcommand{\RatioLatOneFpThirtyTwo}{10.2}
\newcommand{\StudentLatOneBf}{30.2}
\newcommand{\TeacherLatOneBf}{125.6}
\newcommand{\RatioLatOneBf}{4.2}
\newcommand{\StudentMem}{1.96}
\newcommand{\TeacherMem}{25.09}
\newcommand{\RatioMem}{12.8}
\newcommand{\ScrWit}{23.61}
\newcommand{\ScrWitDelta}{+0.00}
\newcommand{\ScrCtlGt}{23.65}
\newcommand{\ScrCtlGtDelta}{+0.04}
\newcommand{\ScrCtlGts}{23.59}
\newcommand{\ScrCtlGtsDelta}{$-$0.02}
\newcommand{\ScrLHalf}{24.20}
\newcommand{\ScrLHalfDelta}{+0.59}
\newcommand{\ScrLOne}{24.86}
\newcommand{\ScrLOneDelta}{+1.25}
\newcommand{\ScrLTwo}{25.44}
\newcommand{\ScrLTwoDelta}{+1.83}
\newcommand{\ScrLFour}{25.53}
\newcommand{\ScrLFourDelta}{+1.92}
\newcommand{\ScrLEight}{25.95}
\newcommand{\ScrLEightDelta}{+2.33}
\newcommand{\ScrLSixteen}{25.88}
\newcommand{\ScrLSixteenDelta}{+2.26}
\newcommand{\ScrSemTwoAll}{26.74}
\newcommand{\ScrSemTwoAllDelta}{+3.13}
\newcommand{\ScrSemTwoChg}{26.70}
\newcommand{\ScrSemTwoChgDelta}{+3.09}
\newcommand{\ScrSemFourAll}{26.40}
\newcommand{\ScrSemFourAllDelta}{+2.79}
\newcommand{\ScrSemFourChg}{26.46}
\newcommand{\ScrSemFourChgDelta}{+2.85}
\newcommand{\ScrFeatA}{26.76}
\newcommand{\ScrFeatADelta}{+3.15}
\newcommand{\ScrFeatB}{26.86}
\newcommand{\ScrFeatBDelta}{+3.25}
\newcommand{\ScrFeatC}{26.29}
\newcommand{\ScrFeatCDelta}{+2.68}
\newcommand{\IncSemTwoAll}{+0.80}
\newcommand{\IncSemTwoChg}{+0.75}
\newcommand{\IncSemFourAll}{+0.45}
\newcommand{\IncSemFourChg}{+0.51}
\newcommand{\IncLSixteen}{$-$0.07}
\newcommand{\IncFeatA}{+0.01}
\newcommand{\IncFeatB}{+0.12}
\newcommand{\IncFeatC}{$-$0.45}
\newcommand{\LkSlots}{4770}
\newcommand{\LkDistinct}{333}
\newcommand{\LkLocations}{12}
\newcommand{\LkTest}{477/477}
\newcommand{\LkVal}{477/477}
\newcommand{\LkPixels}{198.8}
\newcommand{\LkOracleSek}{77.87}
\newcommand{\LkOracleFscd}{95.04}
\newcommand{\LkOracleMiou}{94.25}
\newcommand{\LkKnown}{90.8}
\newcommand{\RndTwoFifty}{59.85}
\newcommand{\RndSixHundred}{69.42}
\newcommand{\RndSixHundredSd}{0.17}
\newcommand{\PerascdLandsat}{65.21}
\newcommand{\RndLastSlope}{0.18}
\newcommand{\FZeroTeacher}{28.61}
\newcommand{\FZeroBase}{28.75}
\newcommand{\FZeroBaseSd}{0.37}
\newcommand{\FZeroKd}{31.21}
\newcommand{\FZeroKdSd}{0.09}
\newcommand{\FZeroGain}{2.46}
\newcommand{\FZeroKdMinusT}{+2.60}
\newcommand{\FZeroTMinusB}{$-$0.14}
\newcommand{\FZeroLocs}{00, 02, 20}
\newcommand{\FOneTeacher}{26.02}
\newcommand{\FOneBase}{25.57}
\newcommand{\FOneBaseSd}{0.08}
\newcommand{\FOneKd}{28.05}
\newcommand{\FOneKdSd}{0.15}
\newcommand{\FOneGain}{2.48}
\newcommand{\FOneKdMinusT}{+2.03}
\newcommand{\FOneTMinusB}{+0.45}
\newcommand{\FOneLocs}{01, 21, 30}
\newcommand{\FTwoTeacher}{22.48}
\newcommand{\FTwoBase}{21.48}
\newcommand{\FTwoBaseSd}{0.32}
\newcommand{\FTwoKd}{23.32}
\newcommand{\FTwoKdSd}{0.08}
\newcommand{\FTwoGain}{1.84}
\newcommand{\FTwoKdMinusT}{+0.84}
\newcommand{\FTwoTMinusB}{+1.00}
\newcommand{\FTwoLocs}{10, 12, 22}
\newcommand{\FThreeTeacher}{20.18}
\newcommand{\FThreeBase}{18.42}
\newcommand{\FThreeBaseSd}{0.60}
\newcommand{\FThreeKd}{21.70}
\newcommand{\FThreeKdSd}{0.17}
\newcommand{\FThreeGain}{3.28}
\newcommand{\FThreeKdMinusT}{+1.52}
\newcommand{\FThreeTMinusB}{+1.75}
\newcommand{\FThreeLocs}{11, 31, 32}
\newcommand{\LPrimGain}{2.53}
\newcommand{\LPrimGainLo}{0.74}
\newcommand{\LPrimGainHi}{4.32}
\newcommand{\LPrimGainP}{=0.03}
\newcommand{\LPrimPos}{9/9}
\newcommand{\LPrimWilcox}{=0.004}
\newcommand{\LPrimFscd}{1.30}
\newcommand{\LPrimMiou}{1.46}
\newcommand{\LPrimGainFinal}{2.48}
\newcommand{\LSecGain}{2.51}
\newcommand{\LSecGainLo}{1.58}
\newcommand{\LSecGainHi}{3.45}
\newcommand{\LSecGainP}{=0.003}
\newcommand{\LSecPos}{12/12}
\newcommand{\LSecWilcox}{=0.0005}
\newcommand{\LSecFscd}{1.22}
\newcommand{\LSecMiou}{1.41}
\newcommand{\LSecGainFinal}{2.51}
\newcommand{\LPrimKdMinusT}{1.46}
\newcommand{\LSecKdMinusT}{1.75}
\newcommand{\LPrimKdMinusTP}{=0.05}
\newcommand{\LSecKdMinusTP}{=0.02}
\newcommand{\LTeacherMinusBMean}{0.77}
\newcommand{\LBaseMean}{23.56}
\newcommand{\LFoldSd}{4.54}
\newcommand{\LKdSlopeMax}{0.10}
\newcommand{\CorrR}{0.09}
\newcommand{\CorrP}{=0.87}
\newcommand{\CorrN}{6}
\newcommand{\TostMargin}{0.25}
\newcommand{\TostP}{=0.02}
\newcommand{\TostLo}{$-$0.01}
\newcommand{\TostHi}{+0.20}
\newcommand{\SecFullMaxFinalGap}{0.02}
\newcommand{\SecFullBestEpMin}{108}
\newcommand{\SecFullBestEpMax}{119}
\newcommand{\SecBaseBestEpMin}{47}
\newcommand{\SecBaseBestEpMax}{100}
\newcommand{\TeacherAdvVitg}{2.21}
\newcommand{\TeacherAdvVmb}{1.41}
\newcommand{\LKdMinusTMaxMin}{0.22}
\newcommand{\ComputeTotal}{479}
\newcommand{\ComputeLandsatKd}{29.0}
\newcommand{\ComputeLandsatTeachers}{25.8}

\renewcommand{\thetable}{S\Roman{table}}
\renewcommand{\thefigure}{S\arabic{figure}}
\renewcommand{\thesection}{S\arabic{section}}

\begin{document}
\title{Supplementary Material for ``DistilSCD: Foundation-Model Accuracy at Lightweight Cost for
Semantic Change Detection through Output Distillation, with a Location-Disjoint Landsat-SCD
Protocol''}
\author{Mathis~Magron, Eric~Hervet, and~Andy~Couturier}
\maketitle
\sloppy

All numbers below are generated from the run logs by the scripts in \texttt{paper/scripts/} of
the code repository; SeK, Fscd and mIoU are in percentage points.

\section{SECOND: Per-Seed Results and the 200-Epoch Extension}
Table~\ref{tab:s-second-runs} lists every SECOND test run. Arms with the same seed share their
initialisation, data order and augmentation draws; their losses at the first optimiser step are
identical, which we checked for every arm. Without distillation, the best epoch falls between
47 and 100 and the SeK then declines under the constant learning rate; with distillation, the
best epoch is at or near the end of the 120-epoch budget. Because the distilled runs were still
improving, we resumed them from epoch 120 to 200 after seeing the curves, a decision we report
as post hoc. The extension adds +\ExtInc\ [\ExtIncLo, \ExtIncHi] points to the best-epoch SeK
(paired, $p\ExtIncP$), reaching \SecFullTwoHundredSek\,$\pm$\,\SecFullTwoHundredSekSd, i.e.
\SecFullTwoHundredVsT\ [\SecFullTwoHundredVsTLo, \SecFullTwoHundredVsTHi] above the teacher
($p\SecFullTwoHundredVsTP$); the curves are flat after epoch 150. Over the same 200 epochs, the
undistilled student's best epoch is unchanged and its final SeK falls to
\SecBaseSekFinalTwoHundred; the distillation gain at 200 epochs is +\SecFullTwoHundredGain\ on
the best epoch and +\SecFullTwoHundredGainFinal\ on the final epoch.

\begin{table}[!t]
\centering
\footnotesize
\caption{SECOND test results per seed.}
\label{tab:s-second-runs}
\begin{tabular}{lrrrrrrr}
\toprule
Arm & Epochs & Seed & SeK max & Best ep. & SeK final & Fscd & mIoU \\
\midrule
change KD (ViT-G) & 120 & 1 & 25.61 & 118 & 25.61 & 65.64 & 74.37 \\
change KD (ViT-G) & 120 & 2 & 25.50 & 115 & 25.47 & 65.53 & 74.33 \\
change KD (ViT-G) & 120 & 3 & 25.41 & 119 & 25.41 & 65.44 & 74.31 \\
change KD (ViT-G) & 120 & 4 & 25.69 & 117 & 25.68 & 65.72 & 74.43 \\
change KD (ViT-G) & 120 & 5 & 25.74 & 114 & 25.70 & 65.77 & 74.42 \\
DistilSCD (ViT-G) & 120 & 1 & 26.39 & 119 & 26.39 & 66.47 & 74.64 \\
DistilSCD (ViT-G) & 120 & 2 & 26.18 & 108 & 26.16 & 66.23 & 74.59 \\
DistilSCD (ViT-G) & 120 & 3 & 26.11 & 119 & 26.11 & 66.23 & 74.52 \\
DistilSCD (ViT-G) & 120 & 4 & 26.14 & 116 & 26.13 & 66.23 & 74.55 \\
DistilSCD (ViT-G) & 120 & 5 & 26.19 & 119 & 26.19 & 66.28 & 74.54 \\
DistilSCD (VMamba-B) & 120 & 1 & 25.69 & 118 & 25.68 & 65.74 & 74.37 \\
DistilSCD (VMamba-B) & 120 & 2 & 25.59 & 119 & 25.59 & 65.62 & 74.34 \\
DistilSCD (VMamba-B) & 120 & 3 & 25.46 & 105 & 25.45 & 65.56 & 74.25 \\
DistilSCD (VMamba-B) & 120 & 4 & 25.90 & 118 & 25.89 & 65.92 & 74.47 \\
DistilSCD (VMamba-B) & 120 & 5 & 25.82 & 119 & 25.82 & 65.78 & 74.46 \\
no KD & 120 & 1 & 23.87 & 100 & 23.73 & 64.45 & 73.18 \\
no KD & 120 & 2 & 23.77 & 94 & 23.61 & 64.35 & 73.14 \\
no KD & 120 & 3 & 23.91 & 62 & 23.48 & 64.45 & 73.20 \\
no KD & 120 & 4 & 23.94 & 66 & 23.55 & 64.39 & 73.31 \\
no KD & 120 & 5 & 23.76 & 47 & 23.37 & 64.19 & 73.25 \\
no KD & 120 & 6 & 24.20 & 62 & 24.04 & 64.72 & 73.36 \\
no KD & 120 & 42 & 23.84 & 72 & 23.73 & 64.38 & 73.18 \\
DistilSCD (ViT-G) & 200 & 1 & 26.45 & 182 & 26.39 & -- & -- \\
DistilSCD (ViT-G) & 200 & 2 & 26.36 & 151 & 26.28 & -- & -- \\
DistilSCD (ViT-G) & 200 & 3 & 26.19 & 143 & 26.08 & -- & -- \\
DistilSCD (ViT-G) & 200 & 4 & 26.25 & 196 & 26.24 & -- & -- \\
DistilSCD (ViT-G) & 200 & 5 & 26.37 & 191 & 26.35 & -- & -- \\
no KD & 200 & 1 & 23.87 & 100 & 23.26 & 64.45 & 73.18 \\
no KD & 200 & 2 & 23.77 & 94 & 23.23 & 64.35 & 73.14 \\
no KD & 200 & 3 & 23.91 & 62 & 23.24 & 64.45 & 73.20 \\
no KD & 200 & 4 & 23.94 & 66 & 23.28 & 64.39 & 73.31 \\
no KD & 200 & 5 & 23.76 & 47 & 23.33 & 64.19 & 73.25 \\
no KD & 200 & 6 & 24.20 & 62 & 23.52 & 64.72 & 73.37 \\
no KD & 200 & 42 & 23.84 & 72 & 23.28 & 64.38 & 73.18 \\
\bottomrule
\end{tabular}
\par\smallskip\parbox{\linewidth}{\scriptsize Fscd and mIoU at the best epoch. The 200-epoch DistilSCD rows resume the 120-epoch runs (post-hoc extension); the 200-epoch no-KD rows are the full runs whose first 120 epochs give the 120-epoch rows.}
\end{table}


\section{Screening Runs}
Table~\ref{tab:s-screening-runs} lists the 32 screening runs (16 arms, two seeds) on the
internal validation split. The split (2\,671 training and 297 validation pairs from the official
training set) is fixed and hashed and is released with the code. The decision rule (``best mean
validation SeK; an increment counts if it exceeds 0.36 points'') was written in the versioned lab journal before each of the three waves (waves 1--3 were trained at commits \texttt{2f8ece0}, \texttt{51205c5} and \texttt{8de169a}, the feature wave at \texttt{945879a}; see documentation/distillation.md). It is a record specified in advance, not a public registration. The change-weight grid was extended twice, from $\{0.5,1,2\}$ to $\{4,8\}$ and then to 16, because the best value lay at its edge. Two caveats apply to the screening only. First, the SECOND teacher was trained
by its authors on all 2\,968 training pairs, including our 297 validation pairs, which can favour
distillation over the ground-truth controls on validation; it cannot explain the gap of about
1.8 points between distillation and the controls at the same weight, but the controls were not run on the test set, so the test confirms the distillation gain, not the control comparison. Second, two seeds detect only large increments; the screening
ranks configurations, and the test confirmation with five seeds supports the conclusions.

\begin{table}[!t]
\centering
\footnotesize
\caption{Screening runs on the SECOND internal validation split (2\,671 / 297 pairs).}
\label{tab:s-screening-runs}
\begin{tabular}{lrrrrrr}
\toprule
Arm tag & Seed & SeK max & Best ep. & SeK final & A100-h & GB \\
\midrule
ctl-gt-l2 & 1 & 23.60 & 76 & 22.87 & 4.7 & 4.2 \\
ctl-gt-l2 & 2 & 23.69 & 51 & 22.68 & 4.7 & 4.2 \\
ctl-gtsmooth-l2 & 1 & 23.60 & 114 & 23.55 & 4.7 & 4.2 \\
ctl-gtsmooth-l2 & 2 & 23.59 & 33 & 22.57 & 4.8 & 4.2 \\
kdchg-l05 & 1 & 23.82 & 119 & 23.82 & 5.0 & 4.2 \\
kdchg-l05 & 2 & 24.58 & 75 & 24.20 & 4.7 & 4.2 \\
kdchg-l1 & 1 & 24.54 & 119 & 24.54 & 4.8 & 4.2 \\
kdchg-l1 & 2 & 25.19 & 91 & 24.78 & 5.0 & 4.2 \\
kdchg-l16 & 1 & 25.88 & 108 & 25.80 & 4.7 & 4.2 \\
kdchg-l16 & 2 & 25.87 & 110 & 25.84 & 4.7 & 4.2 \\
kdchg-l2 & 1 & 25.30 & 118 & 25.29 & 4.6 & 4.2 \\
kdchg-l2 & 2 & 25.59 & 95 & 25.33 & 4.7 & 4.2 \\
kdchg-l4 & 1 & 25.29 & 86 & 25.27 & 4.7 & 4.2 \\
kdchg-l4 & 2 & 25.77 & 82 & 25.70 & 4.7 & 4.2 \\
kdchg-l8 & 1 & 26.23 & 109 & 26.10 & 4.7 & 4.2 \\
kdchg-l8 & 2 & 25.66 & 85 & 25.48 & 4.7 & 4.2 \\
kdfeat-s12-l1 & 1 & 26.87 & 117 & 26.84 & 11.4 & 17.4 \\
kdfeat-s12-l1 & 2 & 26.65 & 85 & 26.49 & 11.0 & 17.4 \\
kdfeat-s12-l4 & 1 & 26.39 & 116 & 26.39 & 10.7 & 17.4 \\
kdfeat-s12-l4 & 2 & 26.19 & 114 & 26.19 & 10.8 & 17.4 \\
kdfeat-s123-l1 & 1 & 26.73 & 119 & 26.73 & 10.7 & 17.4 \\
kdfeat-s123-l1 & 2 & 27.01 & 105 & 26.91 & 10.7 & 17.4 \\
kdsem-l8-t2-all & 1 & 26.79 & 115 & 26.75 & 4.7 & 4.2 \\
kdsem-l8-t2-all & 2 & 26.70 & 105 & 26.57 & 4.6 & 4.2 \\
kdsem-l8-t2-changed & 1 & 26.61 & 114 & 26.60 & 4.7 & 4.2 \\
kdsem-l8-t2-changed & 2 & 26.78 & 91 & 26.78 & 4.7 & 4.2 \\
kdsem-l8-t4-all & 1 & 26.27 & 114 & 26.27 & 4.7 & 4.2 \\
kdsem-l8-t4-all & 2 & 26.52 & 107 & 26.30 & 4.7 & 4.2 \\
kdsem-l8-t4-changed & 1 & 26.68 & 119 & 26.68 & 4.7 & 4.2 \\
kdsem-l8-t4-changed & 2 & 26.25 & 118 & 26.23 & 4.6 & 4.2 \\
witness & 1 & 23.40 & 66 & 23.07 & 4.7 & 4.2 \\
witness & 2 & 23.82 & 93 & 23.44 & 5.0 & 4.2 \\
\bottomrule
\end{tabular}
\par\smallskip\parbox{\linewidth}{\scriptsize Tags: kdchg-l$\lambda$ = change KD with weight $\lambda$; ctl-gt / ctl-gtsmooth = ground-truth controls; kdsem-l8-t$T$-mask = semantic KD; kdfeat-s(stages)-l$\lambda$ = feature KD; witness = no KD.}
\end{table}


\section{Teachers}
\subsection{Reproduction of the Published SECOND Checkpoints}
Table~\ref{tab:s-teacher-repro} reports the two published teachers re-evaluated with our code.
For the ViT-G teacher, the authors' evaluation code and ours agree to $10^{-8}$ on the same
predictions; the deformable-attention CUDA operator and its PyTorch reference agree on
99.9994\,\% of the arg-max decisions; our SECOND copy and the authors' pre-processed copy are
equivalent for evaluation (no pixel where the change map and the semantic labels disagree, among
444\,M pixels, and identical class order). Half precision changes the ViT-G score by less than
0.01 point, so the cache was computed in float16. The VMamba-B checkpoint is sensitive to the
input normalisation: with the PerA statistics used by the released evaluation code it scores
23.73, and with ImageNet statistics it reproduces exactly the score stored in the checkpoint;
its float16 inference is also degraded, so its cache was computed in float32 and stored in
float16. We verified that the selective-scan kernels were not the cause (compiled kernel,
PyTorch reference scan, and TF32 disabled give the same SeK on a 256-pair subset).

\begin{table*}[!t]
\centering
\footnotesize
\caption{Re-evaluation of the published SECOND teachers on the full test set (1\,694 pairs).}
\label{tab:s-teacher-repro}
\begin{tabular}{lllrrrr}
\toprule
Teacher & Setting & Normalisation & SeK & Fscd & mIoU & ms (b8) \\
\midrule
PerASCD ViT-G & precision fp32 & PerA & 26.108 & 66.410 & 74.331 & 1613.4 \\
PerASCD ViT-G & precision fp16 & PerA & 26.108 & 66.411 & 74.331 & 489.7 \\
PerASCD ViT-G & precision bf16 & PerA & 26.105 & 66.409 & 74.330 & 562.6 \\
VMamba-B & precision fp32 & imagenet & 25.314 & 65.614 & 74.014 & 504.6 \\
VMamba-B & precision fp16 & imagenet & 23.659 & 64.419 & 72.956 & 182.4 \\
VMamba-B & precision bf16 & imagenet & 25.312 & 65.615 & 74.013 & 187.4 \\
VMamba-B & normalisation, fp32 & PerA (released evaluation code) & 23.732 & -- & -- & -- \\
VMamba-B & normalisation, fp32 & ImageNet & 25.314 & -- & -- & -- \\
VMamba-B & normalisation, fp32 & dataset statistics (commented in released code) & 24.528 & -- & -- & -- \\
VMamba-B & normalisation, fp32 & per-date 0-255 & 24.469 & -- & -- & -- \\
\bottomrule
\end{tabular}
\par\smallskip\parbox{\linewidth}{\scriptsize Published: ViT-G 26.11 (training log 26.1087); VMamba-B 25.31 (stored in the checkpoint: 0.25314). Latency at batch 8 on one A100 (the VMamba-B job ran on a different node).}
\end{table*}


\subsection{Teacher Equivariance and the D4 Cache}
Between $T(g\cdot x)$ and $g\cdot T(x)$, the ViT-G teacher's change decision differs on
2.2--2.5\,\% of the pixels of the SECOND training set and its class decision on 2.7--3.0\,\% of
the changed pixels. With about 20\,\% of pixels changed, a single-view cache would have given
targets inconsistent with the student's augmentation on roughly a tenth of the changed area.
On its own training set, read back from the half-precision cache, the teacher scores 33.24 SeK
(Fscd 71.50, mIoU 77.77): the soft targets are far from a copy of the labels.

\subsection{Landsat-SCD Fold Teachers}
Table~\ref{tab:s-teachers} summarises the four retrained ViT-G teachers. They follow the authors'
recipe (SGD, learning rate 0.1, momentum 0.9, polynomial decay, 10\,\% warm-up, 50 epochs, batch
4 with 2 accumulation steps, float16 autocast, gradient clipping 1.5, drop-path 0.3, D4 and
photometric augmentation), with two declared differences: the checkpoint is selected on the fold's
validation SeK instead of the test Fscd, and activation checkpointing of the 40 ViT blocks is used
to fit a 40\,GB A100. On CPU with a ViT-B and drop-path 0.3, checkpointing gives gradients
identical to the reference (maximum difference 0.0). The 486 PerA pre-trained tensors are loaded
and checked before training.

\begin{table*}[!t]
\centering
\footnotesize
\caption{Retrained PerASCD ViT-G teachers, one per Landsat-SCD fold (50 epochs, authors' recipe).}
\label{tab:s-teachers}
\begin{tabular}{rlrrrlrrrrr}
\toprule
Fold & Test tiles & Val & Ep. & Gate & Pass & Test & Test final & Test max & A100-h & GB \\
\midrule
0 & 00, 02, 20 & 21.85 & 28 & 22.50 & no & 28.61 & 29.12 & 29.99 & 6.4 & 28.5 \\
1 & 01, 21, 30 & 25.74 & 42 & 23.36 & yes & 26.02 & 25.47 & 26.35 & 6.4 & 21.0 \\
2 & 10, 12, 22 & 16.89 & 31 & 14.90 & yes & 22.48 & 23.00 & 23.10 & 6.7 & 23.5 \\
3 & 11, 31, 32 & 23.25 & 11 & 20.70 & yes & 20.18 & 20.19 & 20.74 & 6.4 & 23.5 \\
\bottomrule
\end{tabular}
\par\smallskip\parbox{\linewidth}{\scriptsize Val/Test: SeK at the best-validation epoch; Gate: mean best validation SeK of the three no-KD students of the fold; Test max is shown for information only (never used for selection).}
\end{table*}


\section{Landsat-SCD Protocol}
\subsection{Overlap Verification}
\url{scripts/verify_landsat_leakage.py} hashes the pixels of every image slot of the
released LandsatSCD512 archive (two images per pair), independently of file names. It reports
the number of distinct images, the pairs of each split whose two images exist identically in the
training set, the consistency of the class labels across duplicates, the hypothesis that the tile
index of the name identifies a location (every (index, year) group holds a single distinct image,
and no image is shared by two indices in the same year), and the SeK of a lookup oracle that
labels each test pixel with the classes annotated for the same image in the training set.
Table~II of the paper reports the results. The same hash check finds no exact duplicate between
the SECOND training (5\,936 images) and test (3\,388 images) sets; partial spatial overlap without
exact duplicates is not tested.

\subsection{Fold Construction and Budget}
The folds (\texttt{splits/LandsatSCD\_loc/fold\{0..3\}}) result from an exhaustive search over
the 15\,400 partitions of the 12 tiles into four groups of three, scored on the balance of the nine
transition types and of the change ratio; in each fold, one non-test tile serves as validation.
Table~\ref{tab:s-pilot} reports the budget pilot on fold 0: the rule fixed before the run, read on
validation only, selects the window of epochs 50--99, hence the 100-epoch budget used for every
arm and fold.

\begin{table}[!t]
\centering
\footnotesize
\caption{Budget pilot on Landsat-SCD fold 0 (student without KD, one seed).}
\label{tab:s-pilot}
\begin{tabular}{lrrrr}
\toprule
Epochs & Val mean & Val slope & Test mean & Test slope \\
\midrule
0-49 & 18.56 & 3.576 & 25.19 & 3.772 \\
50-99 & 22.84 & 0.044 & 29.25 & -0.005 \\
100-149 & 22.63 & -0.050 & 29.12 & 0.006 \\
150-199 & 22.49 & -0.007 & 29.12 & -0.023 \\
200-249 & 22.58 & -0.008 & 29.01 & -0.046 \\
250-299 & 22.49 & -0.011 & 28.77 & -0.024 \\
300-349 & 22.36 & -0.032 & 28.76 & -0.004 \\
350-399 & 22.16 & -0.071 & 28.71 & -0.004 \\
400-449 & 22.11 & -0.004 & 28.71 & -0.002 \\
\bottomrule
\end{tabular}
\par\smallskip\parbox{\linewidth}{\scriptsize Slopes in SeK points per 10 epochs. Rule fixed before the run, applied to validation only: first 50-epoch window with a validation slope below 0.05, i.e. epochs 50--99, hence a 100-epoch budget.}
\end{table}


\subsection{All Runs}
Table~\ref{tab:s-landsat-runs} lists the 24 Landsat-SCD runs. Fold-0 distilled runs read the cache
\url{perascd_landsat_f0_train_exploratory}, written only because the gate was lifted
explicitly; the lift is recorded in the cache directory and in each run's configuration.

\begin{table*}[!t]
\centering
\footnotesize
\caption{Landsat-SCD location-disjoint folds, all runs.}
\label{tab:s-landsat-runs}
\begin{tabular}{rlrrrrrrrr}
\toprule
Fold & Arm & Seed & Val max & Best ep. & Test & Fscd & mIoU & Test final & Val slope \\
\midrule
0 & no KD & 1 & 22.77 & 61 & 29.07 & 72.07 & 75.73 & 28.88 & 0.009 \\
0 & no KD & 2 & 22.47 & 97 & 28.35 & 71.60 & 75.33 & 28.34 & 0.109 \\
0 & no KD & 3 & 22.27 & 99 & 28.84 & 71.99 & 75.58 & 28.84 & 0.089 \\
0 & DistilSCD & 1 & 24.41 & 77 & 31.19 & 72.83 & 76.81 & 31.24 & -0.044 \\
0 & DistilSCD & 2 & 24.08 & 82 & 31.31 & 72.99 & 76.84 & 31.45 & 0.020 \\
0 & DistilSCD & 3 & 24.68 & 99 & 31.13 & 72.78 & 76.84 & 31.13 & 0.064 \\
1 & no KD & 1 & 23.44 & 98 & 25.58 & 69.43 & 74.57 & 25.59 & 0.113 \\
1 & no KD & 2 & 23.39 & 44 & 25.65 & 69.47 & 74.61 & 25.79 & -0.053 \\
1 & no KD & 3 & 23.25 & 35 & 25.49 & 69.29 & 74.53 & 25.76 & 0.043 \\
1 & DistilSCD & 1 & 25.72 & 70 & 27.93 & 70.14 & 75.97 & 28.11 & -0.035 \\
1 & DistilSCD & 2 & 25.52 & 99 & 28.00 & 70.25 & 76.04 & 28.00 & 0.021 \\
1 & DistilSCD & 3 & 25.53 & 98 & 28.22 & 70.43 & 76.13 & 28.22 & 0.030 \\
2 & no KD & 1 & 15.27 & 74 & 21.74 & 68.52 & 73.09 & 21.61 & -0.052 \\
2 & no KD & 2 & 14.56 & 98 & 21.12 & 68.03 & 72.72 & 21.11 & 0.143 \\
2 & no KD & 3 & 14.87 & 84 & 21.58 & 68.46 & 72.99 & 21.61 & 0.012 \\
2 & DistilSCD & 1 & 18.73 & 93 & 23.41 & 69.23 & 73.95 & 23.39 & 0.097 \\
2 & DistilSCD & 2 & 17.52 & 69 & 23.28 & 69.06 & 73.88 & 23.38 & 0.023 \\
2 & DistilSCD & 3 & 17.49 & 90 & 23.27 & 69.11 & 73.89 & 23.31 & 0.079 \\
3 & no KD & 1 & 20.84 & 60 & 18.23 & 63.88 & 71.18 & 18.57 & -0.024 \\
3 & no KD & 2 & 21.06 & 89 & 19.09 & 64.65 & 71.65 & 19.10 & 0.091 \\
3 & no KD & 3 & 20.21 & 57 & 17.94 & 63.64 & 71.09 & 18.06 & -0.008 \\
3 & DistilSCD & 1 & 23.97 & 97 & 21.56 & 66.32 & 73.15 & 21.57 & 0.027 \\
3 & DistilSCD & 2 & 24.14 & 97 & 21.89 & 66.26 & 73.36 & 21.87 & 0.038 \\
3 & DistilSCD & 3 & 23.55 & 98 & 21.64 & 66.24 & 73.16 & 21.64 & 0.032 \\
\bottomrule
\end{tabular}
\par\smallskip\parbox{\linewidth}{\scriptsize Test, Fscd, mIoU at the best-validation epoch; val slope in SeK points per 10 epochs over epochs 70--99. Fold 0 is exploratory (teacher gate lifted).}
\end{table*}


\section{Student Design Before This Study}
The student was designed earlier by the first author as an efficient SCD model and its recipe was
frozen before the distillation study. Table~\ref{tab:s-student-design} summarises the design
results that fixed the recipe, for transparency. Three blocks inherited from that design (the
checkerboard state-space fusion, the frequency branch and the change-guided attention) have no
detectable effect; the fusion block was removed, the other two remain in the network that is
distilled here. These design experiments selected epochs on the SECOND test set, as is customary
on this dataset; their effects were required to hold on both the best and the final epoch.

\begin{table*}[!t]
\centering
\footnotesize
\caption{Student design results obtained before this study (SECOND, SeK points).}
\label{tab:s-student-design}
\begin{tabular}{p{0.5\linewidth}rrl}
\toprule
Component & Contribution & 95\,\% CI & Verdict \\
\midrule
C2S2 block (checkerboard + MCA-SF + S6 scan), removed & 0.16 & [-0.04, 0.36] & not detectable \\
MCA-SF alone & 0.15 & [-0.24, 0.52] & not detectable \\
change-guided attention (CGA), kept & -0.01 & -- & not detectable \\
FFT frequency branch, kept & -0.46 & -- & not established \\
wider decoder (+3.5 M, +12 GMACs) & -0.07 & -- & not detectable \\
removing the SeK loss term of Mamba-FCS & 1.25 & -- & established \\
temporal-swap augmentation & 0.96 & -- & established \\
EMA of weights & 0.74 & -- & established (0.74-0.87) \\
DySample instead of bilinear upsampling & 0.75 & -- & established \\
\bottomrule
\end{tabular}
\par\smallskip\parbox{\linewidth}{\scriptsize From the design campaign of the student (about 150 runs, 4--8 seeds per configuration, SECOND test with best-epoch and final-epoch readings). These results fixed the student recipe that is then frozen in this paper; they are reported for transparency and are not part of the distillation evidence.}
\end{table*}


\section{Compute}
Table~\ref{tab:s-compute} lists the GPU time used by this study, about \ComputeTotal\ A100-hours
in total, excluding teacher caches and checks (about 10 A100-hours).

\begin{table*}[!t]
\centering
\footnotesize
\caption{Compute used in this study.}
\label{tab:s-compute}
\begin{tabular}{p{0.62\linewidth}rrl}
\toprule
Study & Runs & A100-h & Source \\
\midrule
SECOND: KD screening, 3 logit waves (26 runs) & 26 & 126.8 & sacct \\
SECOND: feature-KD screening (6 runs) & 6 & 66.4 & sacct \\
SECOND: test confirmation, ViT-G teacher (10 runs) & 10 & 66.0 & sacct \\
SECOND: end-to-end trial & 1 & 0.5 & sacct \\
SECOND: VMamba-B teacher study (5 runs + teacher eval/cache) & 5 & 32.6 & sacct \\
SECOND: 200-epoch extension (5 resumes) & 5 & 21.7 & sacct \\
Landsat random split: student baseline, 600 epochs (5 runs, incl. first 250) & 5 & 70.5 & timing.csv \\
Landsat fold 0: 600-epoch budget pilot (stopped at epoch 477) & 1 & 11.2 & timing.csv \\
Landsat folds: student baselines (12 runs) & 12 & 28.2 & timing.csv \\
Landsat folds: PerASCD ViT-G teachers (4 runs) & 4 & 25.8 & timing.csv \\
Landsat folds: KD students (12 runs) & 12 & 29.0 & timing.csv \\
TOTAL (excluding teacher caches and checks, about 10 A100-h) & 87 & 478.7 & sum \\
\bottomrule
\end{tabular}
\end{table*}


\section{Reproducibility}
The repository contains the training and evaluation code (\texttt{distilscd/}), the scripts for
teacher evaluation and caching (\texttt{scripts/teacher/}), the SLURM launchers used on the
Narval cluster (\url{scripts/train_kd_second.sbatch}, \url{scripts/train_kd_landsat.sbatch},
\url{scripts/teacher/teacher_landsat.sbatch}), the splits (\texttt{splits/}), the per-run logs
used for every table, and the scripts that regenerate the tables, numbers and figures of the paper
from those logs (\url{paper/scripts/build_tables.py}, \url{make_numbers.py},
\url{make_tables_tex.py}, \url{make_figures.py}). Each run directory stores the full
configuration and the commit it was trained with; a launcher refuses to reuse a run directory for
a different configuration. GPU training is not bit-reproducible, so replays differ by up to
0.5 SeK on Landsat-SCD; the path without distillation was verified to be bit-identical on CPU
before and after the distillation code was added.

\end{document}
